\documentclass{article}
\usepackage{graphicx} % Required for inserting images
\usepackage{amsmath}
\usepackage{amssymb}
\usepackage{amsthm}
\usepackage{xcolor} % Required for defining custom colors

\definecolor{darkblue}{rgb}{0.0, 0.0, 0.5} % Define a darker blue color

\usepackage[colorlinks=true, linkcolor=darkblue, urlcolor=darkblue, citecolor=darkblue]{hyperref} % Use the darker blue color

\title{Compactness}
\author{Aaron Honculada}
\date{February 2025}

\begin{document}

\maketitle

\section{Why is compactness important?}
\subsection{Generalization of finite sets}
Compactness can be though of as a generalization of \textbf{finite sets} to infinite spaces. Many properties
that hold for finite sets extend to compact spaces:
\begin{itemize}
    \item In finite sets, any sequence has a convergent subsequence.
    \item In compact spaces, any sequence has a \textbf{convergent subsequence} (Bolzano-Weierstrass Theorem).
    \item In finite sets, every function is bounded.
    \item In compact spaces, every \textbf{continuous function} is bounded (Extreme Value Theorem).
\end{itemize}

\subsection{Compact sets behave well under Continuous Maps}
If $f: X \to Y$ is \textbf{continuous} and $X$ is compact, then:
\begin{itemize}
    \item The \textbf{image} $f(X)$ is compact.
    \item If $Y$ is \href{run:../hausdorff/hausdorff.pdf#nameddest=hausdorff-definition}{Hausdorff}, then $f$ is \textbf{closed}.
    \item If $f$ is also \textbf{injective}, then $f$ is \textbf{homeomorphism onto its image}, which
    is very useful in functional analysis.
\end{itemize}

\subsection{Extreme Value Theorem: Maximum and Minimum always exist}
If $f: X \to \mathbb{R}$ is \textbf{continuous} and $X$ is compact, then $f$ attains its maximum and minimum values.
This is crucial in optimization problems, physics, and economics.
Example:
\begin{itemize}
    \item If $f(x)$ is continuous on a closed interval $[a, b]$, then it has a maximum and a minimum, 
    ensuring stability in real-world problems.
\end{itemize}

\subsection{Compactness leads to Convergence: Bolzano-Weierstrass Theorem \& Sequential Compactness}
\begin{itemize}
    \item In compact spaces, every sequence has a \textbf{convergent subsequence}.
    \item This is why compactness is essential in analysis and metric spaces--it guarantees that ``infinite''
    sequences do not escape to infinity.
\end{itemize}
Example:
\begin{itemize}
    \item The \textbf{unit interval} $[0, 1]$ is compact, so any bounded sequence has a convergent subsequence.
    \item The open interval $(0, 1)$ is \textbf{not} compact--sequences can approach 0 or 1, but not converge
    inside of $(0, 1)$.
\end{itemize}


\subsection{Compactness helps prove existence theorems}
Compactness often helps prove that solutions to problems exist.
Example:
\begin{itemize}
    \item In \textbf{differential equations}, compactness is used in the \textbf{Arzel\`{a}-Ascoli Theorem} to prove
    the existence of solutions.
    \item In \textbf{functional analysis}, the \textbf{Riesz Representation Theorem} and \textbf{Hahn-Banach Theorem}
    rely on compactness.
\end{itemize}

\subsection{Heine-Borel Theorem: Compactness in Euclidean Space}
In $\mathbb{R}^n$, a set is compact \textbf{if and only if it is closed and bounded}. This makes 
compactness easy to check in Euclidean spaces.
Example:
\begin{itemize}
    \item The \textbf{closed interval} $[0, 1]$ is compact.
    \item The \textbf{open interval} $(0, 1)$ is \textbf{not} compact, because it is not closed.
\end{itemize}

\subsection{Compactness Topology and Geometry}
\begin{itemize}
    \item \textbf{Compact spaces are ``small" in a certain sense.}
    \begin{itemize}
        \item Example: The unit sphere $S^n$ is compact, while $\mathbb{R}^n$ is not.
    \end{itemize}
    \item Many important theorems in topology hold only for compact spaces.
    \begin{itemize}
        \item The \textbf{Tychonoff Theorem}, which states that any product of compact spaces is compact.
        \item The \textbf{Alexander Duality}, which is a key result in algebraic topology.
    \end{itemize}
\end{itemize}


\subsection{Compactification: Adding Points to Make Spaces Compact}
Sometimes, a space is not compact, but we can \textbf{compactify} it by adding points.
Example:
\begin{itemize}
    \item The \textbf{one-point compactification} of $\mathbb{R}$ adds a point at infinity,
        turning $\mathbb{R}$ into a \textbf{compact space} homeomorphic to $S^1/ \{N\}$.
    \item The \textbf{Stone-Cech Compactification} extends this idea in a more general way.
\end{itemize}


\end{document}